%=====================================================================
% The Evaluation of Finite Bit Sequences
% Build: xelatex -> bibtex -> xelatex -> xelatex
% Bibliography: references.bib (biblatex, backend=bibtex)
% English edition of the German original; the data set underlying
% Section 5 is contained in the repository under dev/random.
%=====================================================================
\documentclass[11pt,a4paper]{scrartcl}

%-----------------------------------------------------
% Fonts (XeLaTeX). Fall back to default fonts if the
% preferred fonts are not installed.
%-----------------------------------------------------
\usepackage{fontspec}
\IfFontExistsTF{EB Garamond}{\setmainfont[Ligatures=TeX]{EB Garamond}}{}
\IfFontExistsTF{Latin Modern Sans}{\setsansfont[Ligatures=TeX]{Latin Modern Sans}}{}
\IfFontExistsTF{Latin Modern Mono}{\setmonofont[Ligatures=TeX]{Latin Modern Mono}}{}

%-----------------------------------------------------
% Packages
%-----------------------------------------------------
\usepackage[left=3cm,right=3.5cm,top=3cm,bottom=3.5cm]{geometry}
\usepackage[english]{babel}
\usepackage{amsmath,amssymb,amsthm}
\usepackage{graphicx}
\usepackage{booktabs}
\usepackage[normalem]{ulem}   % for \xout in Table 2
\usepackage{xcolor}
\usepackage{listings}
\usepackage{csquotes}
\usepackage{microtype}
\usepackage{pgfplots}
\pgfplotsset{compat=1.18}
\usepackage[backend=bibtex,style=numeric,maxbibnames=9]{biblatex}
\addbibresource{references.bib}
\usepackage[hidelinks,pdfusetitle]{hyperref}

%-----------------------------------------------------
% Pseudocode listings
%-----------------------------------------------------
\lstset{%
  basicstyle=\ttfamily\small,
  frame=single,
  captionpos=b,
  breaklines=true,
  keepspaces=true,
  tabsize=2,
  showstringspaces=false
}
\renewcommand{\lstlistingname}{Algorithm}

%-----------------------------------------------------
% Theorems (shared counter throughout)
%-----------------------------------------------------
\theoremstyle{plain}
\newtheorem{satz}{Theorem}
\newtheorem{lemma}[satz]{Lemma}
\newtheorem{korollar}[satz]{Corollary}
\theoremstyle{definition}
\newtheorem{defi}[satz]{Definition}
\newtheorem{beispiel}[satz]{Example}
\theoremstyle{remark}
\newtheorem{bem}[satz]{Remark}

%-----------------------------------------------------
% Title
%-----------------------------------------------------
\subject{A one-step test and the constructive generation of bit
  sequences with exponentially cumulated frequency distributions of
  their runs}
\title{The Evaluation of Finite Bit Sequences}
\author{Nicolas Windt}
\date{October 28, 2026}

\begin{document}

\maketitle

%-----------------------------------------------------
% Abstract
%-----------------------------------------------------
\begin{abstract}
We introduce the \emph{one-step test}, an evaluation procedure for
finite bit sequences which cumulates the frequencies of the runs and
sub-runs of a sequence. Applying it to the output of the random device
\texttt{/dev/random} of a Linux installation yields an almost exactly
decreasing exponential course of the cumulated frequencies, in
accordance with the geometric distribution of run lengths of
independent bits. Motivated by this observation, we study bit
sequences whose cumulated frequency distribution falls off exactly
exponentially and collect them into classes $B_N$. We show that the
distribution of the runs without sub-runs (base distribution) of the
elements of $B_N$ satisfies a strict recursion, and we derive an
identity for the number of offset sub-runs. Building on this, we give
an iterative and a recursive algorithm for generating an element of
$B_N$. Finally, we assign to every natural number a unique
composition of pairwise distinct classes $B_N$, obtain from it a
family of composite frequency distributions, and discuss the notion of
randomness thereby induced.
\end{abstract}

\vspace{0.5em}
\noindent\textbf{Keywords:}
random sequence $\cdot$ run $\cdot$ runs test $\cdot$ cumulated
frequency distribution $\cdot$ exponential distribution $\cdot$
random number generator $\cdot$ \texttt{/dev/random}

%-----------------------------------------------------
% 1  Introduction
%-----------------------------------------------------
\section{Introduction}\label{sec:einleitung}

The starting point of this work are Karl Popper's investigations of
the concept of probability in \enquote{The Logic of Scientific
Discovery} \cite{popper1935}: randomness is not a property that can be
read off a single event, but an ascription tied to the repeatability
and verifiability of observations. For the finite bit sequences
treated here we draw the consequence not to presuppose randomness
axiomatically, but to examine finite sequences by means of
observable, purely combinatorial structures -- without recourse to
limit statements about infinite sequences.

In this work, first an evaluation procedure for finite bit sequences
is designed, the one-step test (Section~\ref{sec:einschritttest}), and
the result of applying it to the output of a random device is analysed
(Section~\ref{sec:experiment}). The experiment will show that the
frequency distribution of the runs and sub-runs of the output of a
random generator approaches an exponentially shaped distribution.
Motivated by this experimental result, the structure of the frequency
distributions of such bit sequences is investigated formally, thereby
explaining the experimental observation (Section~\ref{sec:klassen}).
Finally, algorithms generating elements of the classes arising there
are given (Section~\ref{sec:algorithmen}), and it is shown how these
classes compose via the binary representation of natural numbers
(Section~\ref{sec:zusammengesetzt}). Section~\ref{sec:diskussion}
discusses the consequences for the concept of randomness;
Section~\ref{sec:schluss} summarises.

The contributions in detail:
\begin{itemize}
  \item a linear evaluation procedure which cumulates the frequencies
        of runs \emph{and} sub-runs of finite bit sequences (one-step
        test, Section~\ref{sec:einschritttest});
  \item an experimental evaluation of the output of \texttt{/dev/random}
        (Section~\ref{sec:experiment});
  \item the classes $B_N$ of bit sequences with exactly exponentially
        decreasing cumulated frequency distribution, including the
        structure of their base distribution
        (Theorem~\ref{satz:grundverteilung}) and an offsetting identity
        for sub-runs (Lemma~\ref{lem:abgeltungsidentitaet});
  \item an iterative and a recursive generator for elements of $B_N$
        (Section~\ref{sec:algorithmen});
  \item a one-to-one correspondence between natural numbers and
        compositions of pairwise distinct classes $B_N$
        (Section~\ref{sec:zusammengesetzt}).
\end{itemize}

%-----------------------------------------------------
% 2  State of research
%-----------------------------------------------------
\section{State of Research}\label{sec:stand}

The mathematical foundation of the concept of randomness is the
Kolmogorov axiomatisation \cite{kolmogoroff1933}; the strong and the
weak law of large numbers following from it justify the expectation
that in sufficiently long sequences of independent, uniformly
distributed bits the run lengths are geometrically distributed: the
probability of a run length $k$ is $2^{-k}$ \cite[and the exercises on
runs therein]{feller1968}. The exponentially decreasing course of the
(cumulated) frequencies observed here is therefore to be expected for
ideal random sequences; what is new in this work is not the decrease
itself but the cumulation over the sub-runs and the study of those
finite sequences whose cumulated distribution falls off
\emph{exactly} exponentially.

Statistical procedures for testing random sequences have long been
established. The classical runs test of Wald and Wolfowitz
\cite{waldwolfowitz1940} tests the number of runs of a sequence against
its expected distribution; to this day it is part of the standardised
test batteries for cryptographic random generators, such as the NIST
suite \cite{rukhin2010nist} and its predecessors
\cite{knuth1997,marsaglia1996}. Surveys on the design and testing of
random number generators are given by Knuth \cite{knuth1997} and
L'Ecuyer \cite{lecuyer1994}. The algorithmic concept of randomness of
Martin-Löf \cite{martinloef1966} defines randomness of infinite
sequences via the impossibility of effective tests; closest to our
approach is the idea there of capturing randomness by passing
effective tests, with the difference that we consider finite
sequences and a single, purely combinatorial quantity.

As experimental source we use the device \texttt{/dev/random} of a
Linux installation; its way of operating is documented and has been
analysed repeatedly \cite{gutterman2004}. The evaluation used here is
independent of the particular source; we single out
\texttt{/dev/random} merely as a notoriously good supplier of bit
sequences whose origin is not to be explained causally.

The epistemological framing follows Popper \cite{popper1935}. We
return to the question whether, and in which sense, the classes $B_N$
can replace the probabilistic concept of randomness in
Section~\ref{sec:diskussion}.

%-----------------------------------------------------
% 3  Foundations: runs and sub-runs
%-----------------------------------------------------
\section{Foundations: Runs and Sub-Runs}\label{sec:grundlagen}

\begin{defi}[Set of finite bit sequences]\label{defi:bitfolgen}
\begin{equation}\label{eqn:mengederbitfolgen}
  \mathcal B := \left\{(b_m)_{1\leq m\leq M<\infty},\ b_m\in\left\{0,1\right\}\right\}
\end{equation}
\end{defi}

\noindent We start from an arbitrary finite bit sequence
\begin{equation}\label{eqn:beliebigebitfolge}
  (b_m)_{1\leq m\leq M,\ M<\infty}\in\mathcal B
\end{equation}
For example
\begin{equation}\label{eqn:runtestfolge}
  10001010100111110000011111110001010101
\end{equation}
This bit sequence can also be viewed as a sequence of runs, where a
run is a sequence of constant bits. In our example these are
\begin{equation}\label{eqn:runtestfolgevonlaeufen}
  1\cdot 000\cdot 1\cdot 0\cdot 1\cdot 0\cdot 1\cdot 00\cdot 11111\cdot
  00000\cdot 1111111\cdot 000\cdot 1\cdot 0\cdot 1\cdot 0\cdot 1\cdot 0\cdot 1
\end{equation}

\begin{defi}[Run]\label{defi:laeufe}
A \emph{run} of a bit sequence $(b_m)_{1\leq m\leq M}\in\mathcal B$ is a
constant sub-sequence $(b_p,\dots,b_q)$, maximal in both directions,
with $1\leq p\leq q\leq M$:
\begin{equation}\label{eqn:mengederlaeufe}
  b_p=b_{p+1}=\dots=b_q,\qquad
  \left(p=1\ \text{or}\ b_{p-1}\neq b_p\right),\qquad
  \left(q=M\ \text{or}\ b_q\neq b_{q+1}\right).
\end{equation}
The number $q-p+1$ is called the \emph{length} of the run.
\end{defi}

\begin{lemma}[Run decomposition]\label{lem:lzerlegung}
Every bit sequence $(b_m)_{1\leq m\leq M}\in\mathcal B$ decomposes
uniquely into runs. With $m_1:=1$ and the recursive rule
\begin{equation}\label{eqn:zerlegung}
  m_{i+1}:=\min\left(\left\{m>m_i:\ b_m\neq b_{m_i}\right\}\cup\{M+1\}\right)
\end{equation}
one obtains unique decomposition points
$1=m_1<m_2<\dots<m_N<m_{N+1}=M+1$ with runs
$(b_{m_i},\dots,b_{m_{i+1}-1})$ and run lengths $r_n:=m_{n+1}-m_n$.
In particular
\begin{equation}
  \sum_{n=1}^{N} r_n=M,
\end{equation}
and consecutive runs carry complementary bit values. The sequence
$(r_n)_{1\leq n\leq N}\in\mathbb N$ is called the \emph{run-length
sequence} of $(b_m)$.
\end{lemma}

\begin{proof}
Since $(b_m)$ is constant on each block $[m_i,m_{i+1})$ and
$b_{m_{i+1}}\neq b_{m_i}$ holds (provided $m_{i+1}\leq M$),
(\ref{eqn:zerlegung}) is well defined and yields a strictly
increasing sequence which reaches $M+1$ after finitely many steps.
The maximality conditions of (\ref{eqn:mengederlaeufe}) ensure that
each block is a run in the sense of the definition. Uniqueness
follows because the first run must begin at $m=1$ and every
subsequent run is determined by the end of its predecessor.
\end{proof}

\begin{bem}[Complement and run-length sequence]\label{bem:komplement}
Two bit sequences possess the same run-length sequence if and only if
they are equal or binary complements of each other: together with the
bit of the first run, the run-length sequence determines the bit
sequence uniquely because of the alternating construction, and
passing to the complement changes only that single bit.
\end{bem}

\noindent Since only the lengths of the runs enter all frequency
considerations of this work, we henceforth identify runs with their
length and work with run-length sequences $(r_n)\in\mathbb N$.

Furthermore the trivial observation holds that a run of length $N$
contains runs $s<N$ with the corresponding frequency distribution
\begin{equation}\label{eqn:haufigkeitenteillaeufe}
  _NS=(_Ns_n)_{0<n\leq N},\quad\text{with}\ _Ns_n=N-n+1.
\end{equation}
For illustration we take the run $N=8$
\begin{equation}\label{eqn:teillaufindemlauf}
  1\;1\;1\;1\;1\;1\;1\;1\quad\text{or}\quad 0\;0\;0\;0\;0\;0\;0\;0
\end{equation}
Then $N$ contains:
\begin{align*}
  {}_8s_1&=8-1+1=8\ \text{runs } 1, & {}_8s_2&=8-2+1=7\ \text{runs } 2,\\
  {}_8s_3&=8-3+1=6\ \text{runs } 3, & {}_8s_4&=8-4+1=5\ \text{runs } 4,\\
  {}_8s_5&=8-5+1=4\ \text{runs } 5, & {}_8s_6&=8-6+1=3\ \text{runs } 6,\\
  {}_8s_7&=8-7+1=2\ \text{runs } 7, & {}_8s_8&=8-8+1=1\ \text{run } 8.
\end{align*}

\begin{defi}[Sub-run]\label{defi:teillauf}
The runs of smaller length contained in a run are henceforth called
sub-runs.
\end{defi}

%-----------------------------------------------------
% 4  The one-step test
%-----------------------------------------------------
\section{The One-Step Test}\label{sec:einschritttest}

We view the bit sequence under investigation as a sequence of runs
\begin{equation}\label{eqn:laeuflangefolge}
  (r_n)_{0<n\leq N},\quad r_n\in\mathcal \mathbb N,\quad N\text{ suitable}
\end{equation}
To evaluate such a sequence we determine the cumulated distribution
\begin{equation}\label{eqn:verteilunglaeufeteillaeufe}
  L=(l_p)_{0<p\leq P<\infty}
\end{equation}
of the frequencies $l_p$ of the runs and of the sub-runs contained in
them of the sequence $(r_n)_{0<n\leq N<\infty}$.

We motivate this approach by the following consideration: for the
output $(b_m)\in \mathcal B$ of a random algorithm delivering
independent, uniformly distributed bits, one has
$P(b_{n-1}=b_n)=P(b_{n-1}\neq b_n)=\tfrac{1}{2}$ \cite{feller1968}.
Instead of resorting to limit statements, we prefer to consider, for
each bit $b_m$ of $(b_m)$, the frequency of the cases
$(b_{m})=(b_{m+i})$. This leads us to the approach of cumulated
distributions: we view $(b_m)\in\mathcal B$ as a sequence $(r_n)$ and
determine the frequencies of its runs and of the sub-runs contained
in them.

Trivially, the largest index $P$ of the sequence
(\ref{eqn:verteilunglaeufeteillaeufe}) is given by the longest run of
$(r_n)$:
\begin{equation}\label{eqn:maximalelaengevonl}
  P=\max\left\{ r_1,\dots, r_N\right\}
\end{equation}

\noindent The preceding considerations lead to the following
procedure, to which we give the name one-step test.

\begin{figure}[htbp]
\begin{lstlisting}[caption={One-step test for sequences of runs},label={algo:einschritt-laeufe}]
Input: an arbitrary sequence of runs of length N, N finite
create an array L of length at most N
for each run of the sequence
    determine the length of the run (here l)
    for each position i = 1 to l
        increase the frequency at position i by l - i + 1
output L
\end{lstlisting}
\end{figure}

\noindent The same procedure can be applied to the bit sequence
$(b_m)$ or $(\overline b_m)$ underlying the sequence of runs $(r_n)$.
A slightly different computational rule results; that it yields the
same result as Algorithm~\ref{algo:einschritt-laeufe} is shown in
Lemma~\ref{lem:aequivalenz}.
\begin{figure}[htbp]
\begin{lstlisting}[caption={One-step test for bit sequences},label={algo:einschritt-bitfolge}]
Input: an arbitrary finite sequence of bits of length P, P finite
create an array L of length at most P
create a counter z with value 0
for each bit of the sequence (step by step)
    has the value of the bit changed compared to the previous step?
    if no, increase the value of counter z by 1
    if yes, reset counter z to 1
    for each position i = 1 to z    (for every single bit!)
        increase the frequency at position i by 1
output L
\end{lstlisting}
\end{figure}
The name \enquote{one-step test} derives from the fact that this
version of the algorithm -- the chronologically first -- processes the
input step by step; its inner loop runs for every single bit with the
current counter value $z$.

\begin{lemma}[Equivalence]\label{lem:aequivalenz}
Algorithms \ref{algo:einschritt-laeufe} and
\ref{algo:einschritt-bitfolge} yield the same cumulated distribution
for every bit sequence.
\end{lemma}

\begin{proof}
Consider a run of length $l$. Through this run,
Algorithm~\ref{algo:einschritt-laeufe} increases the frequency at
position $i$ ($1\leq i\leq l$) by $l-i+1$.
Algorithm~\ref{algo:einschritt-bitfolge} has the counter value $z=j$
at the $j$-th bit of the run; this bit therefore increases positions
$1$ through $j$. Position $i$ is thus increased by the bits
$j=i,\dots,l$, that is, exactly $l-i+1$ times. Both algorithms enter
the same contribution for each run, and the contributions of
different runs add up independently.
\end{proof}

%-----------------------------------------------------
% 5  Experiment
%-----------------------------------------------------
\section{A Motivating Experiment}\label{sec:experiment}

We resort to an experimental means in order to produce a bit sequence
whose structure is not to be explained causally: a so-called random
sequence. Such an experimental means is the device
\texttt{/dev/random} of an arbitrary Linux installation
\cite{gutterman2004}, which enjoys the reputation of being a
notoriously good supplier of such bit sequences, even though it works
rather slowly. Moreover, its way of operating can be verified
without gaps. We collected the output of \texttt{/dev/random} over a
period of about fifteen months (November 13, 2016, to February 15,
2018) with a background script which hourly takes exactly one block
of up to $1024$ bytes,
\begin{quote}
\verb|dd if=/dev/random of="$(date +%Y_%m_%d_%H_%M_%S)" bs=1024 count=1|
\end{quote}
(Short read accesses of the entropy pool occasionally result in
smaller files.) The result are $9\,856$ files totalling
$9\,793\,146$ bytes.\footnote{By the definition of the one-step test,
$l_1$ counts every bit of the sequence exactly once; the consistency
condition $l_1=M=78\,345\,168$ is checked automatically during the
evaluation. The data set as well as the collection and evaluation
scripts are contained in the repository (\texttt{dev/random} and
\texttt{paper/DE\_de/experiment}).}
The result of this experimental procedure is a finite bit sequence
whose length $M$ is determined by the experiment:
\begin{equation}\label{eqn:bitfolgebeliebigerlaenge}
  (b_m)_{1\leq m\leq M<\infty},\quad b_m\in\{0,1\},
\end{equation}
with $M=78\,345\,168$ bits.

\noindent As set out in Section~\ref{sec:grundlagen} on runs and
sub-runs, we now view the sequence $(b_m)$ of
(\ref{eqn:bitfolgebeliebigerlaenge}) as a sequence of the form
(\ref{eqn:laeuflangefolge}) and evaluate it with the one-step test.
The result of the evaluation of the output of the device
\texttt{/dev/random} is recorded in Table~\ref{tab:haeufigkeiten}.
The distribution of the runs and sub-runs of the output of
\texttt{/dev/random} falls off exponentially: $\log_2 l_n$ decreases
by almost exactly $1$ per step as $n$ grows. We present these results
graphically in Figure~\ref{fig:haeufigkeiten}; the theoretical series
shown there and in Table~\ref{tab:haeufigkeiten} is the expectation
$l_n=M\cdot 2^{\,1-n}$ for independent, uniformly distributed bits. A
chi-square goodness-of-fit test of the distribution of the run
lengths against the geometric distribution yields $\chi^2=19.8$ at
$22$ degrees of freedom (critical value at level $5\%$: $33.9$); the
estimate of the transition probability
$\hat p=0.49994\pm 0.00006$ lies $1.05$ standard deviations below
$\tfrac12$. The survey is thus compatible with independent, uniformly
distributed bits. In the spirit of the intersubjective verification of
this experiment, the reader is expressly encouraged to reproduce it by
the means of his choice.

\begin{table}[htbp]
\centering
\begin{tabular}{rrrrrrrr}
\toprule
length & frequency & $\log_2$ & expectation & length & frequency & $\log_2$ & expectation\\
\midrule
$1$&$78345168$&$26.22$&$78345168$&$13$&$19282$&$14.23$&$19127$\\
$2$&$39177249$&$25.22$&$39172584$&$14$&$9588$&$13.23$&$9564$\\
$3$&$19592313$&$24.22$&$19586292$&$15$&$4753$&$12.21$&$4782$\\
$4$&$9798825$&$23.22$&$9793146$&$16$&$2364$&$11.21$&$2391$\\
$5$&$4900914$&$22.22$&$4896573$&$17$&$1174$&$10.20$&$1195$\\
$6$&$2451773$&$21.23$&$2448287$&$18$&$561$&$9.13$&$598$\\
$7$&$1227083$&$20.23$&$1224143$&$19$&$248$&$7.95$&$299$\\
$8$&$614018$&$19.23$&$612072$&$20$&$110$&$6.78$&$149$\\
$9$&$307704$&$18.23$&$306036$&$21$&$50$&$5.64$&$75$\\
$10$&$153948$&$17.23$&$153018$&$22$&$22$&$4.46$&$37$\\
$11$&$76935$&$16.23$&$76509$&$23$&$11$&$3.46$&$19$\\
$12$&$38572$&$15.24$&$38254$&$24$&$5$&$2.32$&$9$\\
\bottomrule
\end{tabular}
\caption{Frequency $l_n$ of the runs and sub-runs of the output of
\texttt{/dev/random}, computed from the complete data set
($M=78\,345\,168$ bits in $9\,856$ files; the longest run has length
$25$). The third column gives $\log_2 l_n$, the fourth the expectation
$M\cdot 2^{\,1-n}$ (rounded) for independent, uniformly distributed
bits.}
\label{tab:haeufigkeiten}
\end{table}

\begin{figure}[htbp]
\centering
\begin{tikzpicture}
\begin{axis}[
  width=0.95\textwidth,
  height=8.5cm,
  xlabel={length $n$ of the (sub-)run},
  ylabel={$\log_2$ of the frequency $l_n$},
  xmin=0, xmax=26.5,
  grid=major,
  legend pos=north east,
  legend cell align=left,
]
\addplot[only marks, mark=*, mark size=1.4pt]
  table[x=n, y=logl] {haeufigkeiten.dat};
\addplot[dashed, thick]
  table[x=n, y=theo] {haeufigkeiten.dat};
\legend{measured frequency $l_n$,
        expectation $M\cdot 2^{\,1-n}$}
\end{axis}
\end{tikzpicture}
\caption{Logarithmic representation of the results of
Table~\ref{tab:haeufigkeiten}: $\log_2 l_n$ versus $n$ (points) with
expectation $M\cdot 2^{\,1-n}$ (dashed).}
\label{fig:haeufigkeiten}
\end{figure}

This observation suggests that the output of \texttt{/dev/random} --
or of any device serving the purpose of delivering bit sequences
without causal generation -- exhibits an almost exponentially
decreasing course of the distribution of its runs and contained
sub-runs. Contradictorily, this would explain a structure of chance.
In the following sections we study such bit sequences without regard
to the way they are obtained.

%-----------------------------------------------------
% 6  Bit sequences with exponentially cumulated frequency distributions
%-----------------------------------------------------
\section{Bit Sequences with Exponentially Cumulated Frequency
Distributions}\label{sec:klassen}

\subsection{The cumulated distribution and the classes
\texorpdfstring{$B_N$}{B\_N}}

\begin{defi}[Cumulated distribution]\label{defi:kumulierteverteilung}
Let $(_Nb_m)$ be a bit sequence of $\mathcal B$ and
$(r_n)_{0<n\leq N<\infty}$ its corresponding sequence of runs. We
denote by the sequence
\begin{equation}\label{eqn:defkumuliert}
  _NL:=(\,_Nl_i\,)_{0<i\leq N<\infty}
\end{equation}
the distribution of the frequencies of the runs and contained
sub-runs of the bit sequence $(_Nb_m)$ and call it the cumulated
distribution.
\end{defi}

\begin{defi}[Exponentially decreasing cumulated
distribution]\label{defi:exponentiell}
By the exponentially decreasing cumulated distribution
$_NL:=(\,l_i\,)_{0<i\leq N<\infty}$ of a bit sequence
$(_Nb_m)\in\mathcal B$ we understand
\begin{equation}\label{eqn:exponentiellerverlauf}
  \forall\ 0<n\leq N<\infty:\quad _Nl_n=2^{N-n}
\end{equation}
Equivalent to this is
\begin{align}
  _Nl_N&=1\label{eqn:ln gleich 1}\\
  _Nl_{n-1}&=2\ {_Nl}_n,\quad 2\leq n\leq N\label{eqn:lnminus1 gleich 2 ln}
\end{align}
\end{defi}

\noindent We point out that the cumulated distribution of a bit
sequence does not suffice to characterise the sequence completely.
Rather, this distribution corresponds to a class of bit sequences.
Hence the following definition.

\begin{defi}[Classes $B_N$]\label{defi:bn}
We denote by $B_N$ the set of sequences with longest run $N$ and with
exponentially decreasing cumulated distribution according to
(\ref{eqn:exponentiellerverlauf}) or, equivalently,
(\ref{eqn:ln gleich 1}) and (\ref{eqn:lnminus1 gleich 2 ln}).
\end{defi}

\begin{defi}[Set of classes]\label{defi:bz}
Set of the classes of bit sequences with exponentially decreasing
cumulated distributions:
\begin{equation}\label{eqn:mengederklassen}
  \mathcal{B_Z} := \{B_N,\ 0< N<\infty \}
\end{equation}
\end{defi}

\subsection{The base distribution}

Having fixed the basic definitions, we can turn to the study of the
elements of an arbitrary $B_N$. Starting from the definition of the
exponentially decreasing cumulated distribution, one may investigate
whether the distribution of the frequencies of its runs without
sub-runs follows a systematics. We thus subtract, from each frequency
of the cumulated distribution $_NL$, the already offset sub-runs;
the result we call the base distribution.

\begin{defi}[Base distribution]\label{defi:grundverteilung}
Let $(b_m)\in\mathcal B$ be a bit sequence with run-length sequence
$(r_n)_{1\leq n\leq R}$ and cumulated distribution $L=(l_p)_{1\leq
p\leq P}$. The \emph{base distribution} $G=(g_p)_{1\leq p\leq P}$
counts the runs by length,
\begin{equation}\label{eqn:grunddef}
  g_p:=\left|\left\{n\in\{1,\dots,R\}:\ r_n=p\right\}\right|,\qquad 1\leq p\leq P.
\end{equation}
Cumulated and base distribution are linked by the triangular
relation
\begin{equation}\label{eqn:kumuliertrelation}
  l_p=\sum_{k=p}^{P}\left(k-p+1\right)\cdot g_k,\qquad 1\leq p\leq P,
\end{equation}
for a run of length $k$ contains exactly $k-p+1$ sub-runs of length
$p$. Because of the triangular structure of
(\ref{eqn:kumuliertrelation}), $G$ is uniquely determined by $L$;
the recursive resolution from above,
\begin{equation}\label{eqn:abgeltungrekursiv}
  g_P=l_P,\qquad
  g_p=l_p-\sum_{k=p+1}^{P}\left(k-p+1\right)\cdot g_k
  \quad(p=P-1,\dots,1),
\end{equation}
is precisely the offsetting procedure described in the following.
\end{defi}

For illustration we trace the resolution (\ref{eqn:abgeltungrekursiv})
on an example from $B_5$. In this case the values of
$_5L=(_5l_i)_{0<i\leq5}$ are
\begin{equation}
  _5l_i=2^{5-i}
\end{equation}
that is, explicitly:
\begin{itemize}
  \item $_5l_1=16$,
  \item $_5l_2=8$,
  \item $_5l_3=4$,
  \item $_5l_4=2$,
  \item $_5l_5=1$.
\end{itemize}
The longest run has length $5$ and is not a sub-run. We therefore
subtract the sub-runs offset by the run of length $5$ and obtain:
\begin{itemize}
  \item $5$ sub-runs $1$,
  \item $4$ sub-runs $2$,
  \item $3$ sub-runs $3$,
  \item $2$ sub-runs $4$,
  \item $1$ run $5$.
\end{itemize}
Now the run $5$ is offset, and with it the two runs $4$ as well. The
runs still to be accounted for are:
\begin{itemize}
  \item $_5l_1-5=11$ runs $1$,
  \item $_5l_2-4=4$ runs $2$,
  \item $_5l_3-3=1$ run $3$,
  \item $_5l_4-2=0$ runs $4$,
  \item $_5l_5-1=0$ runs $5$.
\end{itemize}
We now claim the run $3$, since it is no longer present as a
sub-run. This means
\begin{itemize}
  \item $3$ sub-runs $1$,
  \item $2$ sub-runs $2$,
  \item $1$ run $3$.
\end{itemize}
Now the run $3$ is offset. The runs still to be accounted for are:
\begin{itemize}
  \item $_5l_1-5-3=8$ runs $1$,
  \item $_5l_2-4-2=2$ runs $2$,
  \item $_5l_3-3-1=0$ runs $3$.
\end{itemize}
We now claim the runs $2$. This means
\begin{itemize}
  \item $4$ sub-runs $1$,
  \item $2$ sub-runs $2$.
\end{itemize}
Now the runs $2$ are offset. The runs still to be accounted for are:
\begin{itemize}
  \item $_5l_1-5-3-4=4$ runs $1$.
\end{itemize}
Since runs of length $1$ contain no sub-runs, they are offset
immediately. The result of the resolution
(\ref{eqn:abgeltungrekursiv}) is the base distribution
\begin{equation}
  _5G=(\;4;2;1;0;1\;).
\end{equation}
One possible realisation of this distribution is
\begin{equation}\label{eqn:folfgeausb5}
  (\;0101001100011111\;)\ \text{with}\
  (\;\underbrace{ 0\; 1\; 0\; 1}_{4}\;
   \underbrace{00\; 11}_{2}\;
   \underbrace{000}_{1}\;
   \underbrace{\;\;\;\;}_{0}\;
   \underbrace{11111}_{1}\; )\in B_5
\end{equation}
Of course, every other permutation of this sequence of runs is valid;
the bit values are complemented where necessary. By way of example we
repeat this computation for bit sequences of the classes $B_1$
through $B_8$:
\begin{align}
  (\;1\;)_{[B_1]}\label{eqn:folfgeausb1}\\
  (\;0\;;\;1\;)_{[B_2]}\nonumber\\
  (\;1\;;\;0\;;\;1\;)_{[B_3]}\nonumber\\
  (\;2\;;\;1\;;\;0\;;\;1\;)_{[B_4]}\nonumber\\
  (\;4\;;\;2\;;\;1\;;\;0\;;\;1\;)_{[B_5]}\nonumber\\
  (\;8\;;\;4\;;\;2\;;\;1\;;\;0\;;\;1\;)_{[B_6]}\nonumber\\
  (\;16\;;\;8\;;\;4\;;\;2\;;\;1\;;\;0\;;\;1\;)_{[B_7]}\nonumber\\
  (\;32\;;\;16\;;\;8\;;\;4\;;\;2\;;\;1\;;\;0\;;\;1\;)_{[B_8]}
  \label{eqn:folfgeausb8}
\end{align}
We immediately note that the base distributions of the elements of
the classes $B_N$ exhibit a simple structure as $N$ grows: the
frequencies of the runs $2$ through $N$ shift one position to the
right, and the frequency of the runs $1$ doubles. We now underpin this
experimental observation formally.

\begin{satz}\label{satz:grundverteilung}
Let $N\geq 2$ and $(_Nb_m)\in B_N$. Then the base distribution
$(_Ng_n)_{1\leq n\leq N}$ of the runs of $(_Nb_m)$ satisfies
\begin{align}
  _Ng_n&=2^{N-2-n},\quad 1\leq n\leq N-2,\label{eqn:gnform}\\
  _Ng_{N-1}&=0,\label{eqn:g_nminus1}\\
  _Ng_N&=1.\label{eqn:g_n}
\end{align}
In particular, $_NG$ satisfies the recursion
$_Ng_n=2\,{_Ng}_{n+1}$ for $1\leq n\leq N-3$. (For $N=1$ one has
$_1g_1=1$.)
\end{satz}

\begin{proof}
Since the longest run has length $N$, we have $P=N$ and
$_Nl_n=2^{N-n}$ for all $1\leq n\leq N$ (equation
(\ref{eqn:exponentiellerverlauf})). We resolve the recursion
(\ref{eqn:abgeltungrekursiv}) from top to bottom.

For $n=N$ one has $_Ng_N=_Nl_N=2^{0}=1$, which is
(\ref{eqn:g_n}). For $n=N-1$ one has
\begin{equation*}
  _Ng_{N-1}=_Nl_{N-1}-2\cdot _Ng_N=2-2=0,
\end{equation*}
which is (\ref{eqn:g_nminus1}).

For the remaining indices we prove the closed form (\ref{eqn:gnform})
by downward induction over $n=N-2,\dots,1$. Subtracting equations
(\ref{eqn:kumuliertrelation}) for $n$ and $n+1$ yields
\begin{equation}\label{eqn:differenz}
  {}_Nl_n-{}_Nl_{n+1}= {}_Ng_n+\sum_{k=n+1}^{N} {}_Ng_k,\qquad 1\leq n\leq N-1,
\end{equation}
for the contribution of a run of length $k>n$ to $l_n$ exceeds its
contribution to $l_{n+1}$ by exactly $1$.

\emph{Base case} $n=N-2$: from (\ref{eqn:abgeltungrekursiv}) follows
\begin{equation*}
  _Ng_{N-2}=_Nl_{N-2}-2\cdot\underbrace{_Ng_{N-1}}_{=0}
  -3\cdot\underbrace{_Ng_N}_{=1}
  =4-3=1=2^{N-2-(N-2)}.
\end{equation*}

\emph{Induction step}: let $n\leq N-3$ and suppose
$_Ng_k=2^{N-2-k}$ has been shown for all $n<k\leq N-2$. From
(\ref{eqn:differenz}) and
${}_Nl_n-{}_Nl_{n+1}=2^{N-n-1}$ follows, with the geometric sum,
\begin{equation*}
  _Ng_n
  =2^{N-n-1}-\sum_{k=n+1}^{N-2}2^{N-2-k}
  -\underbrace{_Ng_{N-1}}_{=0}-\underbrace{_Ng_N}_{=1}
  =2^{N-n-1}-\left(2^{N-2-n}-1\right)-1
  =2^{N-2-n}.
\end{equation*}
This shows (\ref{eqn:gnform}) for all $1\leq n\leq N-2$ and hence the
claim.
\end{proof}

\begin{bem}[A peculiarity of the method]\label{bem:keinlaufnminus1}
No element of $B_N$, $N\geq 2$, contains a run of length $N-1$:
$_Ng_{N-1}=0$. The base distribution jumps directly from
$_Ng_{N-2}=1$ to $_Ng_N=1$. This is not an irregularity but a
particular, systematic property of the procedure: in the cumulation,
the frequency share of length $N-1$ has been absorbed entirely by the
sub-runs of the unique run of length $N$ -- $_Nl_{N-1}=2$ counts
precisely its two sub-runs of length $N-1$ -- and in the resolution
from above this share is subtracted before an own run of length $N-1$
could become free.
\end{bem}

\begin{lemma}[Sum identity]\label{lem:summenidentitaet}
For all $m\geq 0$,
\begin{equation}\label{eqn:summenidentitaet}
  T_m:=\sum_{k=1}^{m} k\cdot 2^{m-k}=2^{m+1}-m-2.
\end{equation}
\end{lemma}

\begin{proof}
Induction over $m$. For $m=0$ the sum is empty and
$T_0=0=2^{1}-0-2$. In the step,
\begin{equation*}
  T_{m+1}=\sum_{k=1}^{m+1}k\cdot 2^{m+1-k}
  =2\sum_{k=1}^{m}k\cdot 2^{m-k}+(m+1)
  =2T_m+(m+1)
  =2^{m+2}-(m+1)-2.\qedhere
\end{equation*}
\end{proof}

\begin{korollar}[Converse and existence]\label{kor:umkehrung}
If a bit sequence has the run distribution
\begin{equation}\label{eqn:grundform}
  _Ng_n=2^{N-2-n}\ (1\leq n\leq N-2),\qquad _Ng_{N-1}=0,\qquad _Ng_N=1,
\end{equation}
then it lies in $B_N$. In particular, $B_N$ is non-empty for every
$N\geq 1$.
\end{korollar}

\begin{proof}
We have to show that the cumulated distribution is $_Nl_n=2^{N-n}$.
From (\ref{eqn:kumuliertrelation}) and (\ref{eqn:grundform}) follows
for $1\leq n\leq N-2$, with $m:=N-2-n$,
\begin{align*}
  _Nl_n&=\sum_{k=n}^{N-2}\left(k-n+1\right)\cdot 2^{N-2-k}
  +\left(N-n+1\right)\cdot 1\\
  &=\sum_{j=0}^{m}\left(j+1\right)\cdot 2^{m-j}+\left(m+3\right)
  =\sum_{k=1}^{m+1} k\cdot 2^{m+1-k}+\left(m+3\right)\\
  &=\left(2^{m+2}-m-3\right)+\left(m+3\right)=2^{m+2}=2^{N-n},
\end{align*}
where the penultimate step uses the sum identity
(Lemma~\ref{lem:summenidentitaet}, $m\mapsto m+1$). For $n=N-1$ one
has $_Nl_{N-1}=2\cdot _Ng_N=2=2^{N-(N-1)}$ and for $n=N$ one has
$_Nl_N=_Ng_N=1$. The longest run has length $N$, hence the bit
sequence lies in $B_N$.
\end{proof}

\subsection{The offsetting of sub-runs}

The offsetting of sub-runs, as presented in the procedure for
deriving the base distribution, can also be computed directly: the
frequency $_Nt_n$ of the offset sub-runs of length $n$ is the sum of
the sub-run contributions of all longer runs; a run of length $i$
contains exactly $i-n+1$ sub-runs of length $n$:
\begin{equation}\label{eqn:abgeltung}
  _Nt_n=\sum_{i=n+1}^{N}\left(i-n+1\right)\cdot _Ng_i,\qquad 1\leq n\leq N.
\end{equation}
With $_Nt_n = _Nl_n - _Ng_n$ one obtains the general expression
\begin{equation}\label{eqn:abgegoltete}
  _Nl_n - _Ng_n=\sum_{i=n+1}^{N}\left(i-n+1\right)\cdot _Ng_i,
\end{equation}
which holds for arbitrary base distributions and equivalently
reformulates the triangular relation
(\ref{eqn:kumuliertrelation}). For the exponentially decreasing
cumulated distributions $_Nl_n=2^{N-n}$ and the base distribution of
Theorem~\ref{satz:grundverteilung} we verify three cases.

\noindent$\blacklozenge$ $n=N$: the sum in (\ref{eqn:abgeltung}) is
empty, hence
\begin{equation}
  _Nt_N=0.
\end{equation}

\noindent$\blacklozenge$ $n=N-1$:
\begin{equation}
  _Nt_{N-1}=\left(N-(N-1)+1\right)\cdot\underbrace{_Ng_N}_{=1}=2
  =2^{N-(N-1)}-\underbrace{_Ng_{N-1}}_{=0}.
\end{equation}

\noindent$\blacklozenge$ $1\leq n\leq N-2$: with $_Nl_n=2^{N-n}$ and
$_Ng_n=2^{N-2-n}$ one has
\begin{equation}\label{eqn:tnformel}
  _Nt_n=_Nl_n-_Ng_n=2^{N-n}-2^{N-2-n}=3\cdot 2^{N-2-n}.
\end{equation}
From (\ref{eqn:tnformel}) follows, by an index shift, the shift
structure also visible in the examples (\ref{eqn:folfgeausb1})
through (\ref{eqn:folfgeausb8}): for $1\leq n\leq N-2$,
\begin{equation}\label{eqn:verschiebung}
  _Nt_n=3\cdot 2^{N-2-n}=\,_{N-n+1}t_1,
\end{equation}
for by (\ref{eqn:tnformel}) one has $_Mt_1=3\cdot 2^{M-3}$, so for
$M=N-n+1$ precisely the right-hand side of
(\ref{eqn:verschiebung}). The offset sub-runs of length $n$ of a
class $B_N$ thus coincide with the offset sub-runs of length $1$ of
the class $B_{N-n+1}$.

In summary we record:

\begin{lemma}\label{lem:abgeltungsidentitaet}
Let $N<\infty$. Then
\begin{equation}\label{eqn:abgeltungsidentitaet}
  2^{N} = \left(N+1\right)+\left( 1_{\{N>1\}}\cdot\sum_{n=0}^{N-2} (N-1-n) \cdot 2^n\right)
\end{equation}
\end{lemma}

\begin{proof}
The substitution $k=N-1-n$ transforms the sum into
\begin{equation*}
  \sum_{n=0}^{N-2}\left(N-1-n\right)\cdot 2^{n}
  =\sum_{k=1}^{N-1}k\cdot 2^{N-1-k}=T_{N-1},
\end{equation*}
and by Lemma~\ref{lem:summenidentitaet} ($m=N-1$) one has
$T_{N-1}=2^{N}-N-1$, hence
\begin{equation*}
  \left(N+1\right)+T_{N-1}
  =\left(N+1\right)+\left(2^{N}-N-1\right)=2^{N}.\qedhere
\end{equation*}
\end{proof}

\begin{bem}[Total length]\label{bem:gesamtlaenge}
The right-hand side of (\ref{eqn:abgeltungsidentitaet}) decomposes
the bit length of a representative of $B_{N+1}$ into the
contributions of the individual runs: the unique run of length $N+1$
contributes $N+1$ bits, a run of length $k\leq N-1$ occurs
$2^{N-1-k}$ times and contributes $k$ bits each; length $N$ contributes
nothing, for $_{N+1}g_N=0$
(Remark~\ref{bem:keinlaufnminus1}). The lemma thus states: every
representative of $B_{N+1}$ has exactly $2^{N}$ bits, and every
representative of $B_N$ accordingly exactly $2^{N-1}$ bits.
\end{bem}

%-----------------------------------------------------
% 7  Algorithms
%-----------------------------------------------------
\section{Algorithms for Generating an Element of
\texorpdfstring{$B_N$}{B\_N}}\label{sec:algorithmen}

\noindent Motivated by the results of the preceding section we give
an iterative algorithm computing a class representative of $B_N$ for
$N<\infty$.

\begin{figure}[htbp]
\begin{lstlisting}[caption={Iterative generation of a sequence from B\_N for an arbitrary N},label={algo:iterativ}]
Input: N, the length of the longest run
create an empty bit vector B
append a run of length N
for all i from N-2 down to 1
    append 2^(N-i-2) runs of length i to B
output B
\end{lstlisting}
\end{figure}

\noindent The correctness of the algorithm follows immediately from
Corollary~\ref{kor:umkehrung}: it generates precisely the run
distribution (\ref{eqn:grundform}), and every bit sequence with this
run distribution lies in $B_N$.

\noindent For the recursive algorithm we use the following
construction, which describes the passage from $B_{N-1}$ to $B_N$.

\begin{lemma}[Construction]\label{lem:konstruktion}
Let $N\geq 3$ and $({}_{N-1}b_m)\in B_{N-1}$. Extend every run of
$({}_{N-1}b_m)$ by exactly one bit and insert $2^{N-3}$ runs of
length $1$ at positions where the alternating construction is
preserved (for instance one between each pair of adjacent runs and
one at an end). Then the result lies in $B_N$.
\end{lemma}

\begin{proof}
By Theorem~\ref{satz:grundverteilung}, applied to $B_{N-1}$, the
sequence $({}_{N-1}b_m)$ has
\begin{equation*}
  \sum_{n=1}^{N-3}2^{N-3-n}+1=\left(2^{N-3}-1\right)+1=2^{N-3}
\end{equation*}
runs. After the extension the sequence has run lengths $k+1$ instead
of $k$, as well as $2^{N-3}$ new runs of length $1$. A run of length
$k$ contains $k-n+1$ sub-runs of length $n$; after the extension
therefore $(k+1)-n+1=k-(n-1)+1$. For the cumulated distribution $L'$
of the extended sequence one thus has, for all $2\leq n\leq N$,
\begin{equation*}
  l'_n=\sum_{k\geq n-1}\left(k-(n-1)+1\right)\cdot\,_{N-1}g_k
  =\,_{N-1}l_{n-1}=2^{(N-1)-(n-1)}=2^{N-n},
\end{equation*}
so the extension shifts the cumulated distribution by one position.
The new runs of length $1$ change only $l'_1$: by
Remark~\ref{bem:gesamtlaenge}, $({}_{N-1}b_m)$ has $2^{N-2}$ bits,
the extension adds $2^{N-3}$ bits, and the new runs add another
$2^{N-3}$ bits, hence
\begin{equation*}
  l'_1=2^{N-2}+2^{N-3}+2^{N-3}=2^{N-1}=2^{N-1}.
\end{equation*}
The longest run has length $N$ (extension of the run of length
$N-1$), and $l'_n=2^{N-n}$ holds for all $1\leq n\leq N$; the result
therefore lies in $B_N$.
\end{proof}

\noindent On the example $(0001)\in B_3$: extending all runs yields
$(0000\,11)$ with runs of lengths $4$ and $2$; adding
$2^{4-3}=2$ runs of length $1$, for instance by appending them, gives
$(0000\,11\,0\,1)$, whose base distribution is $(2;1;0;1)$ of $B_4$.

\noindent Lemma~\ref{lem:konstruktion} establishes the correctness of
the following recursive algorithm for computing a class
representative of $B_N$ for $N<\infty$.

\begin{figure}[htbp]
\begin{lstlisting}[caption={Recursive computation of a bit sequence from a B\_N},label={algo:rekursiv}]
Input: an element of B_N with N > 2
repeat for every run
    determine the length of the run and extend it by 1 bit
    if the run had length 1, add 2 runs of length 1
output the result
\end{lstlisting}
\end{figure}

The cases $N=1$ and $N=2$ can be treated separately in an
implementation; a tested reference implementation of all algorithms is
given in Appendix~\ref{sec:code}.

%-----------------------------------------------------
% 8  Composite distributions
%-----------------------------------------------------
\section{Composite Elements from Different
\texorpdfstring{$B_N$}{B\_N}}\label{sec:zusammengesetzt}

Let $B\in\mathbb N$ be arbitrary with its unique binary
representation
\begin{equation}\label{eqn:bindarstellung}
  B = \sum_{n=1}^{N} b_n\cdot 2^{n-1},\quad b_n\in\{0,1\},\ b_N=1,
\end{equation}
where $b_1$ is the least significant bit. By
Remark~\ref{bem:gesamtlaenge} every representative of $B_n$ has
exactly $2^{n-1}$ bits. If we thus assign the class $B_n$ to the bit
$b_n$, the number $B$ corresponds exactly to the composition of the
classes $\{B_n:\ b_n=1\}$, and the composite bit sequence has the
length prescribed by $B$:
\begin{equation}\label{eqn:bsummelaufehaeufigkeit}
  B = \sum_{n=1}^{N} b_n\cdot 2^{n-1}
  = \sum_{\substack{1\leq n\leq N\\ b_n=1}}
  \left|\text{representative of } B_n\right|.
\end{equation}
We can thus identify, in a unique way, an arbitrary natural number
$B$ with a composition of pairwise distinct classes $B_N$ of
$\mathcal{B_Z}$ determined uniquely by this number, and we obtain the
composite frequency distribution $_BH_N$ of the runs of a bit
sequence of the prescribed length $B$. We agree on the following
definition.

\begin{defi}[Random bit sequence]\label{defi:zufaellig}
Every bit sequence $(b_n)$ whose frequency distribution of runs is an
arbitrary composition of frequency distributions of pairwise distinct
classes $B_N \in \mathcal{B_Z}$ we call random.
\end{defi}
\begin{bem}[Classification]\label{bem:einordnung}
This definition is to be read as a proposal: it does not replace any
limit or axiomatic statements but fixes a purely combinatorial
condition for finite sequences. Its relation to the probabilistic
and the algorithmic concepts of randomness is discussed in
Section~\ref{sec:diskussion}.
\end{bem}

The frequency of run length $n$ of an arbitrary composition of
pairwise distinct classes of $\mathcal{B_Z}$ follows from
Theorem~\ref{satz:grundverteilung}: the class $B_n$ contributes (for
$b_n=1$) a single run of length $n$, the class $B_{n+1}$ none (for
$_{n+1}g_n=0$), and every class $B_i$ with $i\geq n+2$ contributes
$2^{\,i-n-2}$ runs of length $n$:
\begin{equation}\label{eqn:hnformel}
  _BH_n=b_n+\sum_{i=n+2}^{N} b_{i}\cdot 2^{\,i-n-2},\qquad 1\leq n\leq N.
\end{equation}
The length of the bit sequence obtained in this way is therefore
\begin{align}
  \sum_{n=1}^{N} n\cdot _BH_n
  &= \sum_{n=1}^{N} n\cdot b_n+\sum_{n=1}^{N}\;\sum_{i=n+2}^{N}
     n\cdot b_{i}\cdot 2^{\,i-n-2}\notag\\
  &= \sum_{i=1}^{N} b_i\left(i+\sum_{n=1}^{i-2} n\cdot 2^{\,i-n-2}\right)
   = \sum_{i=1}^{N} b_i\left(i+T_{i-2}\right)\notag\\
  &= \sum_{i=1}^{N} b_i\cdot 2^{\,i-1}=B,\label{eqn:laengekonsistenz}
\end{align}
using the sum identity (Lemma~\ref{lem:summenidentitaet},
$m=i-2$) in the penultimate step, for $T_{i-2}=2^{i-1}-i$. The
composition is thus consistent with the binary representation
(\ref{eqn:bindarstellung}). We illustrate this state of affairs with
\begin{equation}
  B=101010101_2=341_{10}
\end{equation}
in the following table:
\begin{table}[htbp]
\centering
\begin{tabular}{rcccccccccc}
\toprule
$n$&$b_n$&$1$&$2$&$3$&$4$&$5$&$6$&$7$&$8$&$9$\\\midrule
$1$&$1$&$ 1$&  &&  &&&& \\
$2$&$0$&$\xout    0$&$\xout  1$&   &&&&& \\
$3$&$1$&$  1$&$  0$&$  1$&&&&& \\
$4$&$0$&$\xout  2$&$\xout  1$&$\xout  0$&$\xout 1$&&&& \\
$5$&$1$&$  4$&$  2$&$  1$&$0$&$1$&&& \\
$6$&$0$&$\xout  8$&$\xout   4$&$\xout  2$&$\xout 1$&$\xout 0$&$\xout 1$&& \\
$7$&$1$&$16$&$  8$&$  4$&$2$&$1$&$0$&$1$& \\
$8$&$0$&$\xout {32}$&$\xout {16}$&$\xout   8$&$\xout 4$&$\xout 2$&$\xout 1$&$\xout 0$&$\xout 1$& \\
$9$&$1$&$64$&$32$&$16$&$8$&$4$&$2$&$1$&$0$&$1$\\\midrule
&$\Sigma$&$86$&$42$&$22$&$10$&$6$&$2$&$2$&$0$&$ 1$\\
&$n$&$1$&$2$&$3$&$4$&$5$&$6$&$7$&$8$&$9$\\
&$\Pi$&$86$&$84$&$66$&$40$&$30$&$12$&$14$&$0$&$9$\\
$\Sigma$&$341$&&&&&&&&&\\
\bottomrule
\end{tabular}
\caption{Composite distribution for $B=341$: column-wise
contributions of the classes $B_i$ with $b_i=1$ to the frequency
$H_n$ of the runs of length $n$; row $\Sigma$: $H_n$, row $\Pi$:
$n\cdot H_n$. Struck-out entries belong to classes with $b_i=0$.}
\label{tab:zusammengesetzt}
\end{table}

\noindent We are now able to assign to every natural number a
finite, uniquely determined subset
$\{B_n\}_{n\in\mathbb N} \subset \mathcal{B_Z}$ with $n_1\neq n_2\
\forall n_1,n_2$, and thereby to derive a uniquely determined
composite frequency distribution of runs.

\noindent We now study the logarithmic values of the composite
frequencies. Let
\begin{equation}
  C_n:=\sum_{i=n}^{N} {_B}H_i
\end{equation}
be the cumulated frequency of the runs of length at least $n$ of the
composite sequence, where we set $b_j:=0$ for $j>N$. From
(\ref{eqn:hnformel}) follows, by interchanging the sums, the identity
\begin{equation}\label{eqn:kumulatividentitaet}
  C_n
  = b_n + b_{n+1} + \sum_{j=n+2}^{N} b_j\cdot 2^{\,j-n-1}
  = \frac{B_n+b_n}{2},
\end{equation}
where $B_n:=\sum_{j=n}^{N} b_j\cdot 2^{\,j-n}$ is the natural number
with binary representation $(b_N\,\dots\,b_n)$. Since $b_N=1$, one
has $2^{N-n}\leq B_n\leq 2^{N-n+1}-1$, hence the band statement
\begin{equation}\label{eqn:logband}
  2^{N-n-1}\;\leq\; C_n \;\leq\; 2^{N-n},
  \qquad\text{i.e.}\qquad
  N-n-1\;\leq\;\log_2 C_n\;\leq\;N-n.
\end{equation}
Both bounds are sharp: if only $b_N$ is set, then
$C_n=2^{N-n-1}$ and $\log_2 C_n=N-n-1$ exactly; if all bits are set,
then $C_n=2^{N-n}$. Conversely, if a single bit $b_m$ is changed from
$0$ to $1$, then $C_n$ grows, for $n\leq m-2$, by $1+2^{m-n-1}$ and,
for $n\in\{m-1,m\}$, by $1$; this increase is at most $C_n$ itself
(for $n\leq m-2$ because
$1+2^{m-n-1}\leq 2^{N-n-1}\leq C_n$ with $m\leq N-1$ and $n\leq N-3$,
for $n\in\{m-1,m\}$ because $C_n\geq 2^{N-n-1}\geq 1$). The
logarithm to base two therefore grows by at most $1$. This is the
formal background of the experimental observation of
Section~\ref{sec:experiment}: there the logarithmised frequencies
(third column of Table~\ref{tab:haeufigkeiten}) lie just below the
theoretical values (fourth column).

%-----------------------------------------------------
% 9  Discussion
%-----------------------------------------------------
\section{Discussion}\label{sec:diskussion}

Only the frequency distribution of the runs of a bit sequence decides
its membership in a class $B_N$ with $N<\infty$. The consequence is
that not only the seemingly irregular realisation
(\ref{eqn:folfgeausb5}) lies in $B_5$, but also its most regular
variant $11111\,0\,111\,0\,11\,0\,11\,0$: not the arrangement of the
runs but solely their frequency distribution is decisive. As a
further illustration we may view every bit sequence as a binary number
and compare it with any other bit sequence by the usual order
relations. We take the elements of $B_5$. The largest element
is\footnote{The spaces in the representation serve readability only
and have no bearing on the content.}
\begin{equation}
  11111\ 0\ 111\ 0\ 11\ 0\ 11\ 0
\end{equation}
One recognises immediately that the runs of length greater than one
are sorted in decreasing order and are separated by runs of length
one. The smallest element of $B_5$, on the other hand, is
\begin{equation}
  00000\ 1\ 000\ 1\ 00\ 1\ 00\ 1
\end{equation}
One recognises immediately that it is the binary complement of the
largest element of $B_5$. Furthermore,
\begin{equation}\label{eqn:b5kleinstemitseins}
  1\ 00000\ 1\ 000\ 1\ 00\ 1\ 00
\end{equation}
is the smallest element of $B_5$ beginning with a one, and its binary
complement
\begin{equation}\label{eqn:b5groesstermitnull}
  0\ 11111\ 0\ 111\ 0\ 11\ 0\ 11
\end{equation}
is the largest element beginning with a zero. Moreover, $B_5$
contains no element between (\ref{eqn:b5kleinstemitseins}) and
(\ref{eqn:b5groesstermitnull}).

\noindent A certain structure in the classes $B_N$ can thus be
intuitively surmised. For membership in $B_N$ this recognisability is
irrelevant; only the frequency distribution of the runs decides.

\begin{bem}[Cardinality]\label{bem:kardinalitaet}
The classes $B_N$ are large. An element of $B_N$ ($N\geq 2$) is
determined by the order of its $R=2^{N-2}$ runs and the bit of the
first run (Lemma~\ref{lem:lzerlegung}); hence, with the base
distribution of Theorem~\ref{satz:grundverteilung},
\begin{equation}\label{eqn:kardinalitaet}
  \left|B_N\right|
  = 2\cdot\frac{R!}{\prod_{n=1}^{N} g_n!},
  \qquad R=2^{N-2},\quad g_n=2^{N-2-n}.
\end{equation}
(The formula is confirmed for $N\leq 5$ by complete enumeration.)
By Stirling's formula, (\ref{eqn:kardinalitaet}) implies
\begin{equation}\label{eqn:kardinalitaetasymptotik}
  \log_2\left|B_N\right| = 2^{N-1}-\Theta\!\left(N^2\right).
\end{equation}
The description length of an element of $B_N$ thus grows essentially
with the full bit length $2^{N-1}$ of the sequence. Specifying an
element requires practically the specification of the entire
sequence, whereas verifying membership with the one-step test takes
time linear in the sequence length. The fraction of the class among
all sequences of length $2^{N-1}$ is $2^{-\Theta(N^2)}$ and tends to
$0$ -- a complexity statement going beyond this trivial asymmetry of
specification and verification, for instance a relation to
$P\neq NP$ in the sense of cryptographic assumptions, cannot be
derived from the combinatorics of the $B_N$, however.
\end{bem}

\noindent Whether the classes $\mathcal{B_Z}$ can thereby replace
the probabilistic concept of randomness is a philosophical, not a
combinatorial, question. The definition proposed here
(Definition~\ref{defi:zufaellig}) is deliberately purely
combinatorial and tailored to finite sequences; it stands alongside
the probabilistic concept of randomness \cite{kolmogoroff1933} and
the algorithmic concept of randomness of Martin-Löf
\cite{martinloef1966}, which defines randomness of infinite sequences
by passing all effective tests. Against this background the thesis
may be permitted that it is not the randomness of probability theory
that exists, but the arbitrariness of choosing an element of
$\mathcal{B_Z}$ -- and Remark~\ref{bem:kardinalitaet} shows that this
arbitrariness is measurable: it grows with the full length of the
chosen sequence.

%-----------------------------------------------------
% 10  Summary and outlook
%-----------------------------------------------------
\section{Summary and Outlook}\label{sec:schluss}

We have introduced, with the one-step test, an evaluation procedure
for finite bit sequences which cumulates the frequencies of runs and
sub-runs. The experimental application to the output of
\texttt{/dev/random} ($78.3$ million bits, a collection period of
about fifteen months) shows an exponentially decreasing course, as is
to be expected from the geometric distribution of the run lengths of
independent bits; a chi-square goodness-of-fit test confirms this.
The classes $B_N$ of bit sequences with exactly exponentially
decreasing cumulated distribution possess a strictly recursive base
distribution (Theorem~\ref{satz:grundverteilung}); from this result
an iterative and a recursive generator, and, via the binary
representation of natural numbers, a unique composition of the
classes (Section~\ref{sec:zusammengesetzt}).

A tested reference implementation of all four algorithms including
the membership test for the classes $B_N$ is given in
Appendix~\ref{sec:code}; the equivalence of the two variants of the
one-step test is proved in Lemma~\ref{lem:aequivalenz}, and the
relation to the classical concepts of randomness and the cardinality
statements are discussed in Section~\ref{sec:diskussion}.

Open and subject of further work are:
\begin{itemize}
  \item the extension of the experiment to further sources
        (\texttt{/dev/urandom}, pseudo- and hardware generators) and
        larger samples;
  \item improvements to and maintenance of the present text.
\end{itemize}

%-----------------------------------------------------
% Appendix
%-----------------------------------------------------
\appendix
\section{Reference Implementation}\label{sec:code}

The Python module in Algorithm~\ref{lst:code} implements the
one-step test in both variants (Algorithms
\ref{algo:einschritt-laeufe} and \ref{algo:einschritt-bitfolge}), the
two generators (Algorithms \ref{algo:iterativ} and
\ref{algo:rekursiv}) and the membership test for the classes $B_N$.
The self-test verifies, for $N=1,\dots,16$, that both generators
produce only sequences whose cumulated distribution falls off exactly
exponentially, that the recursive passage leads from $B_N$ to
$B_{N+1}$ (Lemma~\ref{lem:konstruktion}), and that both variants of
the one-step test yield identical results on pseudo-random bits
(Lemma~\ref{lem:aequivalenz}).

\lstinputlisting[language=Python,caption={einschritttest.py: reference implementation},label={lst:code}]{einschritttest.py}

%-----------------------------------------------------
% Bibliography
%-----------------------------------------------------
\printbibliography[title={References}]

\end{document}
